\section{Artin 代数性公理}\label{src:C.7}
\begin{flushright}\itshape
一般而言, 我们对非射影存在性定理的认识极其贫乏, 我希望这终会改变.\\
\normalfont Alexander Grothendieck, 致 Murre 的信, 1962 年 \cite[第 $663$ 页]{Mum10}
\end{flushright}
Artin 代数性公理给出通常可在实际中验证的判别条件, 确保给定的叠是代数叠. Artin 在引入代数叠的同一篇论文 \cite{Art74b} 中证明了这一奠基性结果. 下面给出两个版本: 定理 \bookref{C.7.1} 和 \bookref{C.7.4}. 第一个版本是 Artin 代数化 (\bookref{C.6.8}) 的一个相当简单的推论.
\begin{theorem}[Artin 代数性公理, 第一版]\label{src:C.7.1}
设 $\mathcal X$ 为 $(\Sch/\Bbbk)_{\text{ét}}$ 上的叠. 则 $\mathcal X$ 为 $\Bbbk$ 上局部有限型的代数叠, 当且仅当下列条件成立:
\begin{enumerate}[(1)]
\item (保持极限) 叠 $\mathcal X$ 在 $\Sch/\Bbbk$ 上保持极限, 即对每个 $\Bbbk$-代数系统 $B_\lambda$, 函子
\[
\operatorname{colim}\mathcal X(B_\lambda)\to\mathcal X(\operatorname{colim}B_\lambda)
\]
为范畴等价.
\item (对角的可表性) 对角 $\mathcal X\to\mathcal X\times\mathcal X$ 可表.
\item (万有形式形变的存在性) 每个 $x_0\in\mathcal X(\Bbbk)$ 都具有剩余域为 $\Bbbk$ 的完备 Noether 局部 $\Bbbk$-代数 $(R,\mathfrak m)$ 上的万有形式形变 $\{x_n\}$.
\item (有效性) 对每个剩余域为 $\Bbbk$ 的完备 Noether 局部 $\Bbbk$-代数 $(R,\mathfrak m)$, 自然函子
\[
\mathcal X(\Spec R)\to\varprojlim\mathcal X(\Spec R/\mathfrak m^n)
\]
为范畴等价.
\item (万有性的开性) 对每个来自有限型 $\Bbbk$-概形的态射 $U\to\mathcal X$, 若它在 $u\in U(\Bbbk)$ 处万有 (即形式形变 $\{\Spec\widehat{\mathcal O}_{U,u}/\mathfrak m_u^{n+1}\to\mathcal X\}$ \textit{万有}), 则存在 $u$ 的开邻域 $V$, 使得 $U\to\mathcal X$ 在 $V$ 的每个 $\Bbbk$-点处均\textit{万有}.
\end{enumerate}
\end{theorem}
\begin{proof}
首先注意, 对有限型 $\Bbbk$-概形 $U$ 出发的可表局部有限型态射 $U\to\mathcal X$, 光滑性的无穷小提升判别法 (\bookref{4.7.1}) 蕴含, $U\to\mathcal X$ 光滑, 当且仅当它在所有 $\Bbbk$-点 $u\in U$ 处万有. 对 ($\Rightarrow$), (1) 由习题 \bookref{4.3.33} 成立, (2) 由定理 \bookref{4.2.1} 成立, (4) 由例 \bookref{C.5.2} 成立. 若 $U\to\mathcal X$ 是有限型 $\Bbbk$-概形出发的态射, 则它必然可表且局部有限型. 选择光滑呈示 $U\to\mathcal X$ 和 $x_0$ 的原像 $u\in U(\Bbbk)$, 取形式形变 $\{\Spec\mathcal O_{U,u}/\mathfrak m_u^{n+1}\to\mathcal X\}$, 即得 (3). (5) 由光滑性的开性成立.

反之, 首先注意对角的可表性, 即条件 (2), 蕴含概形 $U$ 出发的每个态射 $U\to\mathcal X$ 都可表, 而保持极限的性质 (1) 蕴含 $U\to\mathcal X$ 局部有限型. 对每个对象 $x_0\in\mathcal X(\Bbbk)$, 我们将构造概形出发的光滑态射 $U\to\mathcal X$ 及 $x_0$ 的原像 $u\in U(\Bbbk)$. 条件 (3)--(4) 保证 $x_0$ 存在有效万有形式形变 $\widehat x:\Spec R\to\mathcal X$, 其中 $(R,\mathfrak m)$ 为剩余域为 $\Bbbk$ 的完备 Noether 局部 $\Bbbk$-代数. 由 Artin 代数化 (\bookref{C.6.8}), 存在有限型 $\Bbbk$-概形 $U$, 点 $u\in U(\Bbbk)$, 态射 $p:U\to\mathcal X$, 同构 $R\cong\widehat{\mathcal O}_{U,u}$ 及相容同构 $p|_{R/\mathfrak m^{n+1}}\xrightarrow\sim\widehat x|_{R/\mathfrak m^{n+1}}$. 由 (5), 可将 $U$ 替换为 $u$ 的开邻域, 使 $U\to\mathcal X$ 在 $U$ 的每个 $\Bbbk$-点处万有 (换言之, 光滑). 另见 \cite[\S$5$]{Art74b}, \cite[推论 $10.11$]{LMB00} 和 \cite[Tag 07Y4]{SP}.
\end{proof}
\begin{remark}\label{src:C.7.2}
实际中, 条件 (1)--(4) 常容易直接验证, 其中 (3) 是 Rim--Schlessinger 判别法 (\bookref{C.4.6}) 的推论, (4) 是 Grothendieck 存在性定理 (\bookref{C.5.3}) 的推论. 还需注意, 有时可将本定理应用于对角 $\mathcal X\to\mathcal X\times\mathcal X$, 即概形 $T$ 上对象 $x,y\in\mathcal X(T)$ 的 Isom 层 $\underline{\Isom}_T(x,y)$, 来建立 (2). 在某些特殊情形中, 可直接验证 (5), 但更常用下一小节说明的良好形变与障碍理论来验证.
\end{remark}
\subsection{通过障碍理论表述 Artin 公理}\label{src:C.7.1-subsection}
我们给出 Artin 代数性公理的一个改进版本, 通常更便于实际验证. 具体而言, 若叠具有良好的形变理论和良好的障碍理论, 则万有性的开性 (\bookref{C.7.1}(5)) 成立.

为表述断言, 需要一些记号. 设 $\xi\in\mathcal X(A)$ 为有限生成 $\Bbbk$-代数 $A$ 上的对象. 设 $M$ 为有限 $A$-模, 记 $A[M]$ 为由 $M^2=0$ 定义的环 $A\oplus M$. 令 $\operatorname{Def}_{\xi}(M)$ 为下述图表的同构类集合:
\[
\begin{tikzcd}
\Spec A\arrow[r,"\xi"]\arrow[d,hook]&\mathcal X\\
\Spec A[M]\arrow[ur,dashed,"\eta"']&
\end{tikzcd}
\]
其中两个延拓 $\eta,\eta':\Spec A[M]\to\mathcal X$ 的同构, 按定义是 $\mathcal X(A[M])$ 中限制到 $\xi$ 为恒等的同构 $\eta\xrightarrow\sim\eta'$. 令 $\operatorname{Aut}_{\xi}(M)$ 为平凡形变 $\xi':\Spec A[M]\to\Spec A\to\mathcal X$ 的自同构群. 注意, 当 $\xi\in\mathcal X(\Bbbk)$ 时, $\operatorname{Def}_{\xi}(\Bbbk)$ 恰是 $\mathcal X$ 在 $\xi$ 处的切空间, 而 $\operatorname{Aut}_{\xi}(\Bbbk)$ 是 $\xi$ 的无穷小自同构群, 即 $\operatorname{Aut}_{\mathcal X(\Bbbk[\epsilon])}(\xi')\to\operatorname{Aut}_{\mathcal X(\Bbbk)}(\xi)$ 的核.
\begin{lemma}\label{src:C.7.3}
设 $\mathcal X$ 为满足强齐性条件 $(\mathrm{RS}_4^*)$ 的 $\Sch/\Bbbk$ 上的预叠. 设 $\xi\in\mathcal X(A)$ 为有限生成 $\Bbbk$-代数 $A$ 上的对象.
\begin{enumerate}[(1)]
\item 对每个 $A$-模 $M$, $\operatorname{Def}_{\xi}(M)$ 与 $\operatorname{Aut}_{\xi}(M)$ 自然为 $A$-模, 且函子
\[
\begin{aligned}
\operatorname{Aut}_{\xi}(-)&:\operatorname{Mod}(A)\to\operatorname{Mod}(A)\\
\operatorname{Def}_{\xi}(-)&:\operatorname{Mod}(A)\to\operatorname{Mod}(A)
\end{aligned}
\]
为 $A$-线性的.
\item 考虑核 $I$ 的平方为零的 $\Bbbk$-代数满射 $B\twoheadrightarrow A$, 令 $\operatorname{Lift}_{\xi}(B)$ 为 $\Spec A\to\Spec B$ 上的态射 $\xi\to\eta$ 的集合, 其中 $(\alpha:\xi\to\eta)\sim(\alpha':\xi\to\eta')$ 当且仅当存在同构 $\beta:\eta\to\eta'$ 满足 $\alpha'=\beta\circ\alpha$. 存在 $\operatorname{Def}_{\xi}(I)$ 在 $\operatorname{Lift}_{\xi}(B)$ 上的作用, 对 $B,I$ 具函子性. 若 $\operatorname{Lift}_{\xi}(B)$ 非空, 则此作用自由且传递.
\end{enumerate}
\end{lemma}
\begin{proof}
可用类似引理 \bookref{C.4.9} 的论证建立此结果. 例如, $x\in A$ 的数乘定义为沿 $A$-代数同态 $A[M]\to A[M]$, $a+m\mapsto a+xm$ 诱导的态射 $\Spec A[M]\to\Spec A[M]$ 拉回. 条件 $(\mathrm{RS}_4^*)$ 蕴含函子 $\mathcal X(A[M\oplus M])\to\mathcal X(A[M])\times_{\mathcal X(A)}\mathcal X(A[M])$ 为等价. 加法 $M\oplus M\to M$ 诱导 $A$-代数同态 $A[M\oplus M]\to A[M]$, 因而诱导函子
\[
\mathcal X(A[M])\times_{\mathcal X(A)}\mathcal X(A[M])\cong\mathcal X(A[M\oplus M])\to\mathcal X(A[M]),
\]
定义 $\operatorname{Def}_{\xi}(M)$ 与 $\operatorname{Aut}_{\xi}(M)$ 上的加法.
\end{proof}
与模问题内在的自同构函子 $\operatorname{Def}_{\xi}(-)$ 和形变函子 $\operatorname{Def}_{\xi}(-)$ 不同, 障碍函子是一项额外数据.
